\DocumentMetadata{
  pdfversion=2.0,
  pdfstandard=ua-2,
  lang=en-US,
  tagging=on,
  tagging-setup={math/setup=mathml-SE,tabsorder=structure}
}
\documentclass[11pt,letterpaper]{article}
\usepackage[margin=0.68in,includefoot]{geometry}
\usepackage{newcomputermodern}
\usepackage{amsmath,amscd,mathtools}
\providecommand{\square}{\mdlgwhtsquare}
\usepackage[hidelinks]{hyperref}
\hypersetup{
  pdftitle={Algebra PhD Exam, May 2018},
  pdfauthor={University of Florida Department of Mathematics},
  pdfsubject={Graduate mathematics examination},
  pdfdisplaydoctitle=true,
  bookmarksopen=true
}
\setcounter{secnumdepth}{0}
\setlength{\parindent}{0pt}
\setlength{\parskip}{0.28em}
\setlength{\emergencystretch}{3em}
\setlength{\leftmargini}{2.15em}
\setlength{\leftmarginii}{2.65em}
\setlength{\leftmarginiii}{3em}
\renewcommand{\labelenumi}{\textbf{\theenumi.}}
\pagestyle{empty}
\raggedbottom
\allowdisplaybreaks
\newenvironment{examproblems}[1]{%
  \begin{enumerate}%
  \setcounter{enumi}{#1}%
  \setlength{\topsep}{0.35em}%
  \setlength{\partopsep}{0pt}%
  \setlength{\itemsep}{0.48em}%
  \setlength{\parsep}{0.18em}%
}{\end{enumerate}}
\newenvironment{examsubparts}{%
  \begin{enumerate}%
  \setlength{\topsep}{0.18em}%
  \setlength{\partopsep}{0pt}%
  \setlength{\itemsep}{0.16em}%
  \setlength{\parsep}{0.1em}%
}{\end{enumerate}}
\newenvironment{examinstructions}{%
  \par\begingroup\itshape
}{%
  \par\endgroup\vspace{0.18em}
}
\makeatletter
\renewcommand\section{\@startsection{section}{1}{0pt}%
  {-1.25ex plus -.25ex minus -.1ex}%
  {0.55ex plus .1ex}%
  {\normalfont\large\bfseries}}
\renewcommand{\@maketitle}{%
  \newpage\null\vskip -1.2em%
  {\footnotesize\bfseries
    \href{https://www.ufl.edu/}{University of Florida}, \href{https://math.ufl.edu/}{Department of Mathematics}\par}%
  \vskip 0.18em%
  \tagstructbegin{tag=Title}%
    {\LARGE\bfseries\href{https://gma.math.ufl.edu/exams/algebra-phd/}{Algebra PhD Exam}, May 2018\par}%
  \tagstructend%
  \vskip 0.35em%
  \tagmcbegin{artifact}%
    \hrule height 1pt%
  \tagmcend%
  \vskip 0.55em%
}
\makeatother
\title{Algebra PhD Exam, May 2018}
\author{}
\date{}
\begin{document}
\maketitle
\thispagestyle{empty}
\begin{examinstructions}
Answer seven problems. Write your answers clearly in complete English sentences. You may quote results (within reason) as long as you state them clearly.
\end{examinstructions}
\begin{examproblems}{0}
\item
Determine the Galois group of the polynomial 
\[
f=x^4-2\in\mathbb{Q}[x].
\]
\item
Let~\(K\) be a field and \(K(x)\) the field of rational functions over~\(K\). Let \(\frac{f}{g}\in K(x)\), with \(\frac{f}{g}\notin K\), where~\(f\) and~\(g\) are relatively prime elements of \(K[x]\).
\begin{examsubparts}
\item[\textnormal{(a)}]
Show that \(\frac{f}{g}\) is transcendental over~\(K\).
\item[\textnormal{(b)}]
Prove that \(K(x)\) is a finite extension of \(K\left(\frac{f}{g}\right)\).
\end{examsubparts}
\item
Let~\(C\) be a category and let~\(\Lambda\) be a set. Suppose that for each \(\lambda\in\Lambda\) we are given an object \(X_\lambda\) in~\(C\).
\begin{examsubparts}
\item[\textnormal{(a)}]
Give the definition of a coproduct of the collection of objects \(\{X_\lambda\}_{\lambda\in\Lambda}\).
\item[\textnormal{(b)}]
Prove that any two coproducts of the \(\{X_\lambda\}_{\lambda\in\Lambda}\) are isomorphic.
\item[\textnormal{(c)}]
Prove that in the category of Abelian groups a coproduct exists for any collection of objects indexed by a set~\(\Lambda\).
\end{examsubparts}
\item
Let~\(F\) be the free group on two generators~\(x\) and~\(y\). Prove that the commutator subgroup is not finitely generated.
\item
Let~\(A\) be a torsion abelian group, and let~\(\mathbb{Q}\) be the (additive) group of the rational numbers.
\begin{examsubparts}
\item[\textnormal{(a)}]
Prove \(A\otimes_{\mathbb{Z}}\mathbb{Q}=\{0\}\).
\item[\textnormal{(b)}]
Prove \(\mathbb{Q}\otimes_{\mathbb{Z}}\mathbb{Q}\simeq\mathbb{Q}\).
\end{examsubparts}
\item
Let~\(R\) be a ring and suppose that we have the following commutative diagram of~\(R\)-modules and module homomorphisms with exact rows. 
{\tagpdfsetup{math/mathml/structelem=false,math/mathml/AF=false,math/alt/use=true}
\[
\begin{CD}
0 @>>> A @>{f}>> B @>{g}>> C @>>> 0 \\
@. @V{\alpha}VV @V{\beta}VV @V{\gamma}VV @. \\
0 @>>> A' @>{f'}>> B' @>{g'}>> C' @>>> 0
\end{CD}
\]
}
\begin{examsubparts}
\item[\textnormal{(a)}]
Show that if~\(\alpha\) and~\(\gamma\) are monomorphisms, then so is~\(\beta\).
\item[\textnormal{(b)}]
Show that if~\(\alpha\) and~\(\gamma\) are epimorphisms, then so is~\(\beta\).
\end{examsubparts}
\item
Let~\(R\) be a commutative ring with \(1\ne 0\). Prove that the set consisting of~\(0\) and all zero divisors in~\(R\) contains at least one prime ideal.
\item
State and prove Hilbert’s Basis Theorem.
\item
Prove that an ideal of a Dedekind domain can be generated by two elements.
\item
This problem has two parts.
\begin{examsubparts}
\item[\textnormal{(a)}]
Prove Schur’s Lemma: If~\(R\) is a ring and~\(M\) a simple~\(R\)-module, then \(D:=\operatorname{Hom}_R(M,M)\) is a division ring.
\item[\textnormal{(b)}]
Give an example of~\(R\) and~\(M\) as above where~\(D\) is not a field.
\end{examsubparts}
\item
Let~\(V\) be a nonzero vector space over a field~\(F\). Prove that \(R=\operatorname{Hom}_F(V,V)\) is a simple ring if and only if~\(V\) is finite-dimensional over~\(F\).
\end{examproblems}
\end{document}
